%==============================================================================
\section{An exploratory prototype: learned cut selection}\label{app:rl}\label{ssec:rl}
%==============================================================================

This appendix records a piece of work carried out during the internship that produced
no result about the SSP. It is kept because it was done and because the negative
outcome is informative about where learning could and could not help, and it is here
rather than in Section~\ref{sec:exact} because nothing in it was measured on this
problem. No claim made anywhere else in this report depends on it.

At a given node the Benders solver has several admissible cuts available: a
combinatorial cut for each candidate order it evaluates, the coverage row, and, with
fractional separation enabled, one cut per violated fractional point. Which to add is
currently decided by a fixed rule, and whether a learned policy would do better is a
natural question. Prof.~Colares raised it during the internship.

A prototype was built, following the cut-selection framework of \citet{Tang2020}. Cut
selection is cast as a Markov decision process in which the state is the current
relaxation solution together with a pool of candidate cuts, each summarised by its
violation, support size, right-hand side and mean coefficient magnitude; an action
selects one candidate; and the reward is the resulting increase in the dual bound. The
policy is a softmax that is linear in those features and is trained by the REINFORCE
policy-gradient rule with a return baseline.

Because the Benders solver requires a commercial solver that was unavailable in the
development environment, the learning components were implemented and trained on a
self-contained cut family rather than on SSP instances. The substitute is knapsack
cover cuts: for a knapsack instance and a minimal cover $C$, the inequality
$\sum_{j\in C}x_j\le\lvert C\rvert-1$ is valid, the separation is elementary, and the
linear relaxations can be solved without a commercial solver. The state, the action,
the reward and the policy are the same objects as in the Benders case; only the feature
extractor and the source of the candidate cuts differ.

\paragraph{What was measured.}
The agent was trained for $2{,}000$ episodes at each of six knapsack sizes, with
learning rate $0.2$ and discount $0.95$, and evaluated on a held-out set of instances
it had not seen. The quantity compared is the mean tightening of the linear bound
achieved over a fixed number of separation rounds, under the learned policy against
uniformly random selection from the same candidate pool. Table~\ref{tab:rl} gives the
result.

\begin{table}[h]\centering
\footnotesize
\begin{tabular}{@{}rrrr@{}}
\toprule
knapsack size $n$ & learned policy & random selection & improvement\\
\midrule
$10$ & $1.3855$ & $0.8734$ & $+58.6\%$\\
$12$ & $1.1639$ & $0.7014$ & $+65.9\%$\\
$20$ & $0.7482$ & $0.5021$ & $+49.0\%$\\
$22$ & $0.6692$ & $0.4568$ & $+46.5\%$\\
$24$ & $0.6383$ & $0.3888$ & $+64.2\%$\\
$26$ & $0.5661$ & $0.3279$ & $+72.7\%$\\
\bottomrule
\end{tabular}
\caption{Mean bound tightening on held-out knapsack instances, learned cut selection
against random cut selection. One training seed per size. The learning machinery works;
the problem is not the SSP.}
\label{tab:rl}
\end{table}

Two things follow, and only two. The learning machinery is correct and does learn: the
policy beats random selection at every size tested, by between $46.5\%$ and $72.7\%$. And
that is a statement about knapsack cover cuts, on a single training seed per size, with
no repetition to separate the effect from training variance. It says nothing about the
SSP. The prototype has not been evaluated on SSP instances, and no claim about its
effect on this problem is made anywhere in this report.

Section~\ref{ssec:further} explains why evaluating it on the SSP is unlikely to be
worthwhile until the bound question of Section~\ref{sec:disc} is resolved: cut
selection chooses among cuts a method can already generate, and
Section~\ref{ssec:disc-locality} measures how little those cuts have left to give.
